% GENERATED FILE. Edit canonical lessons, then run npm run generate:books.
% canonical-source-hash: fnv1a64:e260183ff862c632
% canonical-lessons: SA-C108-palanduh, SA-C108-alukam, SA-C108-vrntakam, SA-C108-kusmandam, SA-C108-narikelam

\chapter{Onion, Potato, Aubergine, Pumpkin, Coconut}
\label{ch:sa-onion-potato-aubergine-pumpkin-coconut}
\hlchaptermodality{108}

\begin{chapteropening}
\textbf{By the end of this chapter:} \emph{I can say an onion, a potato, an aubergine, a pumpkin and a coconut.}

Complete the last lesson of chapter 108: I can say an onion, a potato, an aubergine, a pumpkin and a coconut.
\end{chapteropening}

\section[\texorpdfstring{\sk{पलाण्डुः} (palāṇḍuḥ)}{पलाण्डुः (palāṇḍuḥ)}]{\sk{पलाण्डुः} (palāṇḍuḥ) — an onion}

\label{lesson:SA-C108-palanduh}



\begin{quote}
Before the new one: say the Sanskrit for ginger, then the Sanskrit for garlic.
\end{quote}

\subsection*{You'll want to know: \sk{पलाण्डुः}}
\textbf{\sk{पलाण्डुः}} — \emph{palāṇḍuḥ} — \textquotedblleft{}an onion\textquotedblright{}.

\begin{grammarlens}[title={What you've built}]
One more word for this chapter.
\end{grammarlens}

\subsection*{Your turn}
\begin{itemize}
\raggedright
  \item \emph{Say it:} \emph{palāṇḍuḥ}
  \item \emph{Say it:} \emph{palāṇḍuḥ}, once more
  \item \emph{Say it:} say \emph{palāṇḍuḥ}
  \item \emph{Recall:} say the Sanskrit for ginger, then the Sanskrit for garlic, then say \emph{palāṇḍuḥ} again
\end{itemize}

\begin{tcolorbox}[breakable,colback=teal!4,colframe=teal!35!black,title={Before you move on}]
What does \sk{पलाण्डुः} mean? (\textquotedblleft{}An onion\textquotedblright{}.) Say it once more.
\end{tcolorbox}
\section[\texorpdfstring{\sk{आलुकम्} (ālukam)}{आलुकम् (ālukam)}]{\sk{आलुकम्} (ālukam) — a potato}

\label{lesson:SA-C108-alukam}



\begin{quote}
Before the new one: say the Sanskrit for garlic, then the Sanskrit for an onion.
\end{quote}

\subsection*{You'll want to know: \sk{आलुकम्}}
\textbf{\sk{आलुकम्}} — \emph{ālukam} — \textquotedblleft{}a potato\textquotedblright{}.

\begin{grammarlens}[title={What you've built}]
One more word for this chapter.
\end{grammarlens}

\subsection*{Your turn}
\begin{itemize}
\raggedright
  \item \emph{Say it:} \emph{ālukam}
  \item \emph{Say it:} \emph{ālukam}, once more
  \item \emph{Say it:} say \emph{ālukam}
  \item \emph{Recall:} say the Sanskrit for garlic, then the Sanskrit for an onion, then say \emph{ālukam} again
\end{itemize}

\begin{tcolorbox}[breakable,colback=teal!4,colframe=teal!35!black,title={Before you move on}]
What does \sk{आलुकम्} mean? (\textquotedblleft{}A potato\textquotedblright{}.) Say it once more.
\end{tcolorbox}
\section[\texorpdfstring{\sk{वृन्ताकम्} (vṛntākam)}{वृन्ताकम् (vṛntākam)}]{\sk{वृन्ताकम्} (vṛntākam) — an aubergine}

\label{lesson:SA-C108-vrntakam}



\begin{quote}
Before the new one: say the Sanskrit for an onion, then the Sanskrit for a potato.
\end{quote}

\subsection*{You'll want to know: \sk{वृन्ताकम्}}
\textbf{\sk{वृन्ताकम्}} — \emph{vṛntākam} — \textquotedblleft{}an aubergine\textquotedblright{}.

\begin{grammarlens}[title={What you've built}]
One more word for this chapter.
\end{grammarlens}

\subsection*{Your turn}
\begin{itemize}
\raggedright
  \item \emph{Say it:} \emph{vṛntākam}
  \item \emph{Say it:} \emph{vṛntākam}, once more
  \item \emph{Say it:} say \emph{vṛntākam}
  \item \emph{Recall:} say the Sanskrit for an onion, then the Sanskrit for a potato, then say \emph{vṛntākam} again
\end{itemize}

\begin{tcolorbox}[breakable,colback=teal!4,colframe=teal!35!black,title={Before you move on}]
What does \sk{वृन्ताकम्} mean? (\textquotedblleft{}An aubergine\textquotedblright{}.) Say it once more.
\end{tcolorbox}
\section[\texorpdfstring{\sk{कूष्माण्डम्} (kūṣmāṇḍam)}{कूष्माण्डम् (kūṣmāṇḍam)}]{\sk{कूष्माण्डम्} (kūṣmāṇḍam) — a pumpkin}

\label{lesson:SA-C108-kusmandam}



\begin{quote}
Before the new one: say the Sanskrit for a potato, then the Sanskrit for an aubergine.
\end{quote}

\subsection*{You'll want to know: \sk{कूष्माण्डम्}}
\textbf{\sk{कूष्माण्डम्}} — \emph{kūṣmāṇḍam} — \textquotedblleft{}a pumpkin\textquotedblright{}.

\begin{grammarlens}[title={What you've built}]
One more word for this chapter.
\end{grammarlens}

\subsection*{Your turn}
\begin{itemize}
\raggedright
  \item \emph{Say it:} \emph{kūṣmāṇḍam}
  \item \emph{Say it:} \emph{kūṣmāṇḍam}, once more
  \item \emph{Say it:} say \emph{kūṣmāṇḍam}
  \item \emph{Recall:} say the Sanskrit for a potato, then the Sanskrit for an aubergine, then say \emph{kūṣmāṇḍam} again
\end{itemize}

\begin{tcolorbox}[breakable,colback=teal!4,colframe=teal!35!black,title={Before you move on}]
What does \sk{कूष्माण्डम्} mean? (\textquotedblleft{}A pumpkin\textquotedblright{}.) Say it once more.
\end{tcolorbox}
\section[\texorpdfstring{\sk{नारिकेलम्} (nārikelam)}{नारिकेलम् (nārikelam)}]{\sk{नारिकेलम्} (nārikelam) — a coconut}

\label{lesson:SA-C108-narikelam}



\begin{quote}
Before the new one: say the Sanskrit for an aubergine, then the Sanskrit for a pumpkin.
\end{quote}

\subsection*{You'll want to know: \sk{नारिकेलम्}}
\textbf{\sk{नारिकेलम्}} — \emph{nārikelam} — \textquotedblleft{}a coconut\textquotedblright{}.

\begin{grammarlens}[title={What you've built}]
One more word for this chapter.
\end{grammarlens}

\subsection*{Your turn}
\begin{itemize}
\raggedright
  \item \emph{Say it:} \emph{nārikelam}
  \item \emph{Say it:} \emph{nārikelam}, once more
  \item \emph{Say it:} say \emph{nārikelam}
  \item \emph{Recall:} say the Sanskrit for an aubergine, then the Sanskrit for a pumpkin, then say \emph{nārikelam} again
\end{itemize}

\begin{tcolorbox}[breakable,colback=teal!4,colframe=teal!35!black,title={Before you move on}]
What does \sk{नारिकेलम्} mean? (\textquotedblleft{}A coconut\textquotedblright{}.) Say it once more.
\end{tcolorbox}
